
dispose de l'anneau de Chow $A(X)$ des classes de cycles à équivalence rationelle près [4]. Comme variantes de cet anneau, défini d'ailleurs sous des considérations sensiblement plus générales, on a les anneaux $\operatorname{Gr}^\bullet(X)$ (gradué associé à la $\lambda$-filtration de $K^\bullet(X)$) et $\Gtop^\bullet(X)$ (gradué associé à la filtration de $K^\bullet(X)\simeq K_\bullet(X)$ par la codimension des supports). On a des théories des classes de Chern à valeurs dans chacun de ces trois anneaux (cf. [9] pour le cas de $A(X)$, et ce séminaire pour $\operatorname{Gr}^\bullet(X)$, qui s'envoie dans $\Gtop^\bullet(X)$). On a également une théorie cohomologique ($\ell$-adique, disons, où $\ell$ est premier à la caractéristique de $k$) des classes de Chern (SGA 5 VII)~:
\[
 c_i(\underline{E})\in H^{2i}(X,\underline{Z}_\ell(i)),
\]
(où $\underline{E}$ est un fibré vectoriel sur $X$, ou encore un élément quelconque de $K^\bullet(X)$). Ces théories cohomologiques sont reliées à l'aide des homomorphismes canoniques (compatibles avec les structures multiplicatives)
\begin{equation}\tag{4.1}
\begin{tikzcd}[column sep=large,row sep=large]
 H^{2*}(X,\underline{Z}_\ell(*)) & {} \\
 A(X) \arrow[u,"\gamma"] \arrow[r,"\phi"] & \Gtop^\bullet(X) \\
 {} & \operatorname{Gr}^\bullet(X), \arrow[u,"j"]
\end{tikzcd}
\end{equation}
où on a posé
\[
 H^{2*}(X,\underline{Z}_\ell(*))=\coprod_i H^{2i}(X,\underline{Z}_\ell(i)).
\]
L'homomorphisme $j$ provient du fait que la $\lambda$-filtration de $K^\bullet(X)$ est plus fine que la filtration topologique par la codimension du support (Exp. X), et c'est un isomorphisme modulo torsion (Exp. VII 4.11). L'homomorphisme $\phi$ est defini par exemple dans [RRR], et associe à la classe du cycle de codimension $i$ associé à un sous-schéma fermé $Y$ de $X$ codimension $i$, l'élément $\mathrm{c}\ell^\bullet(\underline{O}_Y)$ de $\operatorname{Filt}^i_{\mathrm{top}}(K^\bullet(X))$ mod $\operatorname{Filt}^{i+1}_{\mathrm{top}}(K^\bullet(X))$. C'est un homomorphisme surjectif, dont nous allons voir encore (4.2) que c'est un isomorphisme modulo torsion. Ainsi, \uline{modulo torsion les trois anneaux} $A(X)$, $\Gtop^\bullet(X)$ \uline{et} $\operatorname{Gr}^\bullet(X)$ \uline{sont canoniquement isomorphes}, donc \uline{à posteriori} les formules de Riemann--Roch exprimées à l'aide de l'un ou l'autre de ses anneaux (tensorisés par $\Q$) sont équivalentes. L'anneau gradué $\operatorname{Gr}^\bullet(X)$, surtout utilisé dans ce séminaire comme substitut à l'anneau de Chow, a l'avantage appréciable d'être défini pour tout topos localement annelé, et en particulier pour tout schéma 
